% Generated from locale/tr/content/first-order-logic/natural-deduction/derivations.tex; do not hand-edit.


\iftag{FOL}
      {\olfileid[tr]{fol}{ntd}{der}}
      {\olfileid[tr]{pl}{ntd}{der}}

\olsection{\usetoken{P}{derivation}}

\begin{explain}
Varsayımın ne olduğunu açıkladık ve çıkarım kurallarını verdik. Doğal
tümdengelimde !!^{derivation}s bunlardan tümevarımsal olarak oluşturulur: her
!!{derivation}, ya tek başına bir varsayımdır ya da ardından doğru bir
çıkarım gelen bir, iki veya üç !!{derivation}s içerir.
\end{explain}

\begin{defn}[!!^{derivation}]
\Article{derivation} \emph{!!{derivation}}, varsayımlar~$\Gamma$'dan
!!a{sentence}~$!A$'ya ulaşan, aşağıdaki koşulları sağlayan ve !!{sentence}s
içeren sonlu bir ağaçtır:
\begin{enumerate}
\item Ağacın en üstündeki !!{sentence}s ya $\Gamma$'dadır ya da ağaçtaki bir
  çıkarım tarafından !!{discharged} olarak işaretlenmiştir.
\item Ağacın en altındaki !!{sentence}, $!A$'dır.
\item Ağacın en altındaki tümce~$!A$ dışında ağaçtaki her !!{sentence},
  o !!{sentence} için sonucu doğrudan altında duran doğru bir çıkarım kuralı
  uygulamasının öncülüdür.
\end{enumerate}
Bu durumda !!{derivation} bakımından $!A$'ya \emph{sonuç}, $\Gamma$'ya da
!!{undischarged} varsayımlar deriz.

$\Gamma$'dan $!A$ için !!a{derivation} varsa $!A$'nın $\Gamma$'dan
\emph{!!{derivable}} olduğunu söyleriz; simgesel olarak $\Gamma \Proves
!A$ yazarız. !!a{derivation} $!A$ için bulunuyor ve buradaki her varsayım
!!{discharged} ise $\Proves !A$ yazarız.
\end{defn}

\begin{ex}
Her varsayım tek başına !!a{derivation} olarak alınabilir. Örneğin $!A$ tek
başına !!a{derivation} olur; aynı şey $!B$ için de geçerlidir. Bunlardan,
örneğin $\Intro{\land}$ kuralını uygulayarak yeni bir !!{derivation} elde
edebiliriz:
\begin{prooftree}
\AxiomC{$!A$}
\AxiomC{$!B$}
\RightLabel{\Intro{\land}}
\BinaryInfC{$!A \land !B$}
\end{prooftree}
Bu kurallar genel olacak şekilde tasarlanmıştır: kuraldaki $!A$ ve~$!B$'yi
herhangi bir !!{sentence}s ile, örneğin $!C$ ve~$!D$ ile değiştirebiliriz.
O zaman sonuç $!C \land !D$ olur; dolayısıyla
\begin{prooftree}
\AxiomC{$!C$}
\AxiomC{$!D$}
\RightLabel{\Intro{\land}}
\BinaryInfC{$!C \land !D$}
\end{prooftree}
doğru bir !!{derivation} olur. Elbette varsayımları yer değiştirip $!D$'yi
$!A$'nın, $!C$'yi de $!B$'nin yerine koyabiliriz. Böylece
\begin{prooftree}
\AxiomC{$!D$}
\AxiomC{$!C$}
\RightLabel{\Intro{\land}}
\BinaryInfC{$!D \land !C$}
\end{prooftree}
de doğru bir !!{derivation} olur.

Şimdi başka bir kuralı, örneğin bir koşul sonucu çıkarmamızı ve o koşulun
önbileşeniyle özdeş herhangi bir varsayım bakımından !!{discharge} olanağı veren
$\Intro{\lif}$ kuralını uygulayabiliriz. Dolayısıyla aşağıdakilerin ikisi de
doğru !!{derivation}s olur:
\begin{prooftree}
\AxiomC{$\Discharge{!C}{1}$}
\AxiomC{$!D$}
\RightLabel{\Intro{\land}}
\BinaryInfC{$!C \land !D$}
\DischargeRule{\Intro{\lif}}{1}
\UnaryInfC{$!C \lif (!C \land !D)$}
\DisplayProof\bottomAlignProof
\AxiomC{$!C$}
\AxiomC{$\Discharge{!D}{1}$}
\RightLabel{\Intro{\land}}
\BinaryInfC{$!C \land !D$}
\DischargeRule{\Intro{\lif}}{1}
\UnaryInfC{$!D \lif (!C \land !D)$}
\end{prooftree}
Bunlar sırasıyla $!D \Proves !C \lif (!C \land !D)$ ve
$!C \Proves !D \lif (!C \land !D)$ olduğunu gösterir.

Varsayımları kapatmanın bir zorunluluk değil, izin olduğunu unutmayın:
varsayımları kapatmak zorunda değiliz. Özellikle, varsayımlar
!!{derivation} içinde bulunmasa bile bir kuralı uygulayabiliriz. Örneğin
aşağıda !!{discharged} olacak bir $!A$ varsayımı bulunmamasına rağmen bu
uygulama kurallara uygundur:
\begin{prooftree}
  \AxiomC{$!B$}
  \DischargeRule{\Intro{\lif}}{1}
  \UnaryInfC{$!A \lif !B$}
\end{prooftree}
\end{ex}

